% QUESTAO
No levantamento das alturas de vinte alunos de uma turma, foram obtidos os seguintes valores (em cm): 174, 158, 179, 178, 158, 168, 176, 165, 140, 164, 169, 188, 177, 176, 168, 154, 142, 185, 158, e 185. Qual histograma melhor representa essa distribuição de frequências?

% ALTERNATIVAS
\parbox{\linewidth}{\centering\tikz[xscale=0.35, yscale=0.2]{\draw[CorSerie!20, step=1] (-0.2,0) grid (5.2,7.5);\foreach \i [count=\j] in {2,4,5,6,3} {\pgfmathsetmacro{\int}{30 + \j*15}\draw[thick, fill=CorSerie!\int] (\j-1,0) rectangle (\j, \i);}\draw[thick, CorSerie] (0,0) -- (5,0);\draw[thick, CorSerie] (0,0) -- (0,7.5);}}
\parbox{\linewidth}{\centering\tikz[xscale=0.35, yscale=0.2]{\draw[CorSerie!20, step=1] (-0.2,0) grid (5.2,7.5);\foreach \i [count=\j] in {2,4,5,3,6} {\pgfmathsetmacro{\int}{30 + \j*15}\draw[thick, fill=CorSerie!\int] (\j-1,0) rectangle (\j, \i);}\draw[thick, CorSerie] (0,0) -- (5,0);\draw[thick, CorSerie] (0,0) -- (0,7.5);}}
\parbox{\linewidth}{\centering\tikz[xscale=0.35, yscale=0.2]{\draw[CorSerie!20, step=1] (-0.2,0) grid (5.2,7.5);\foreach \i [count=\j] in {3,6,4,2,5} {\pgfmathsetmacro{\int}{30 + \j*15}\draw[thick, fill=CorSerie!\int] (\j-1,0) rectangle (\j, \i);}\draw[thick, CorSerie] (0,0) -- (5,0);\draw[thick, CorSerie] (0,0) -- (0,7.5);}}
\parbox{\linewidth}{\centering\tikz[xscale=0.35, yscale=0.2]{\draw[CorSerie!20, step=1] (-0.2,0) grid (5.2,7.5);\foreach \i [count=\j] in {2,5,3,6,4} {\pgfmathsetmacro{\int}{30 + \j*15}\draw[thick, fill=CorSerie!\int] (\j-1,0) rectangle (\j, \i);}\draw[thick, CorSerie] (0,0) -- (5,0);\draw[thick, CorSerie] (0,0) -- (0,7.5);}}
\parbox{\linewidth}{\centering\tikz[xscale=0.35, yscale=0.2]{\draw[CorSerie!20, step=1] (-0.2,0) grid (5.2,7.5);\foreach \i [count=\j] in {2,3,6,4,5} {\pgfmathsetmacro{\int}{30 + \j*15}\draw[thick, fill=CorSerie!\int] (\j-1,0) rectangle (\j, \i);}\draw[thick, CorSerie] (0,0) -- (5,0);\draw[thick, CorSerie] (0,0) -- (0,7.5);}}
